\chapter{Filières de formation}\label{ch:in:education-pipelines}

En l'année se terminant en juin 2024, les universités américaines ont décerné 4\,584 masters codés financial mathematics et 870 codés
financial analytics, un code qui n'existait pas avant 2020. En 2010, le premier code comptait 305 masters. La
même année, l'enquête sur les diplômes décernés recense 25\,543 masters en informatique et 1\,760 doctorats de recherche en physique, et
l'industrie recrute parmi tous ceux-là. Ce chapitre compte les diplômes qui alimentent l'industrie, lit les rapports d'insertion de deux
formations spécialisées pour ce qu'ils disent et ce qu'ils ne disent pas, et examine les doctorats, les concours et
l'apprentissage hors d'un cursus.

\section{Diplômes alimentant chaque type d'entreprise}

\begin{definition}[Classification of Instructional Programs]\label{def:in:education-pipelines:cip}
La \emph{Classification of Instructional Programs}\index{Classification of Instructional Programs} (CIP) est la taxinomie américaine
des domaines d'études, avec des codes à six chiffres et leurs définitions, révisée tous les dix ans~; les établissements déclarent leurs
diplômes à l'enquête fédérale sur les diplômes décernés par code CIP et niveau de diplôme.
\end{definition}

L'enquête fédérale compte des diplômes, et non des personnes ni des programmes~: chaque formation choisit son code. Le code de
financial mathematics, inchangé depuis la classification de 2010, est défini comme «~l'application des mathématiques et de la
statistique à l'industrie financière, y compris le développement, la critique et l'usage de divers modèles financiers~»~; la
révision de 2020 a ajouté financial analytics, «~la modélisation des mégadonnées financières~». Aucun des deux codes ne couvre
toutes les formations en ingénierie financière, et une partie de la croissance du premier peut refléter des formations qui
l'adoptent~; les chiffres constituent donc un plancher pour les diplômes financiers spécialisés, et un guide pour la tendance.

\begin{dated}{2024-06}{Diplômes américains dans les domaines de l'industrie}\label{dat:in:education-pipelines:degrees}
Diplômes décernés en premier cursus au cours de l'année se terminant en juin 2024 (IPEDS)~: financial mathematics 611 licences, 4\,584 masters, 24
doctorats de recherche~; financial analytics 553, 870, 4~; statistique 3\,379, 2\,716, 454~; mathématiques 14\,333, 2\,230,
1\,165~; informatique 44\,627, 25\,543, 1\,628~; physique 6\,008, 1\,934, 1\,760. Masters en financial mathematics~:
305 en 2010, 773 en 2014, 3\,554 en 2019.
\end{dated}

\begin{omfigure}
\begin{tikzpicture}
\begin{axis}[width=10.5cm,height=5.8cm,xmin=2009,xmax=2025,ymode=log,ymin=100,ymax=40000,
  xtick={2010,2014,2019,2024},xticklabel style={/pgf/number format/1000 sep=},xlabel={\small année de diplôme (fin en juin)},
  ylabel={\small masters (échelle logarithmique)},tick label style={font=\footnotesize},grid=major,grid style={gray!20},
  table/col sep=comma,unbounded coords=jump,log basis y=10,
  legend style={font=\scriptsize,at={(0.5,-0.3)},anchor=north,legend columns=3,legend cell align=left,
  /tikz/every even column/.append style={column sep=6pt}}]
\addplot[omThm,thick,mark=square*] table[x=year,y=financial_mathematics]{figdata/industry/29-education-pipelines/masters.csv};
\addplot[omThm!50,thick,mark=triangle*] table[x=year,y=financial_analytics]{figdata/industry/29-education-pipelines/masters.csv};
\addplot[omDef,thick,mark=*] table[x=year,y=statistics]{figdata/industry/29-education-pipelines/masters.csv};
\addplot[black!70,thick,mark=o,mark options={solid}] table[x=year,y=mathematics]{figdata/industry/29-education-pipelines/masters.csv};
\addplot[omDef!50,thick,dashed,mark=diamond*,mark options={solid}] table[x=year,y=computer_science]{figdata/industry/29-education-pipelines/masters.csv};
\addplot[black!40,thick,mark=x] table[x=year,y=physics]{figdata/industry/29-education-pipelines/masters.csv};
\legend{mathématiques financières,analyse financière,statistique,mathématiques,informatique,physique}
\end{axis}
\end{tikzpicture}

\omcaption{Masters décernés aux premiers cursus dans six domaines, années se terminant en juin 2010, 2014, 2019 et 2024
(enquête IPEDS sur les diplômes). L'analyse financière (\emph{financial analytics}) est un code de 2020 et n'apparaît qu'une fois. Données~:
\texttt{data/industry/ipeds\_completions.csv}, via \texttt{in\_edu.completions}.}\label{fig:in:education-pipelines:masters}
\end{omfigure}

Les différentes entreprises recrutent dans des diplômes différents. Les rôles de recherche (chapitre~17) s'appuient sur les doctorats et les masters en
mathématiques, statistique, physique et informatique~; les rôles d'ingénierie (chapitres~19 et~20) s'appuient sur l'informatique à tous les
niveaux~; les rôles quant en banque (chapitre~18) et la structuration (chapitre~23) s'appuient sur les masters financiers spécialisés
ainsi que sur les doctorats. Les diplômes spécialisés sont petits à côté des diplômes généraux (la \cref{fig:in:education-pipelines:masters})~:
les masters de financial mathematics étaient moins d'un cinquième de ceux d'informatique en 2024.

\section{Les masters spécialisés~: taille, coût et insertion}

\begin{definition}[Rapport d'insertion professionnelle]\label{def:in:education-pipelines:placement}
Un \emph{rapport d'insertion professionnelle}\index{rapport d'insertion professionnelle} (placement report) est le compte rendu, établi par la formation elle-même, de ce que ses diplômés ont fait après le
diplôme~: combien ont cherché un emploi, combien ont accepté une offre à une date donnée, où, et le salaire déclaré par ceux qui l'ont
indiqué, selon une norme de déclaration que la formation choisit ou qu'elle n'applique pas.
\end{definition}

Un \omterm{def:in:education-pipelines:placement}{rapport d'insertion professionnelle} est un témoignage d'une partie intéressée, et une formation le reconnaît elle-même~: «~Il n'existe pas de norme commune pour la déclaration des données d'insertion
parmi les masters professionnels en ingénierie financière, et les rapports de ces formations ne font pas l'objet d'un audit.~» Les questions du lecteur portent sur les dénominateurs~: sur combien d'étudiants, combien
ont cherché un emploi, combien ont accepté une offre, et combien
ont déclaré un salaire.

\begin{dated}{2026-01}{Deux rapports d'insertion professionnelle}\label{dat:in:education-pipelines:placement}
UC Berkeley Master of Financial Engineering, promotion 2025~: 70 étudiants en recherche d'emploi, 67 offres acceptées~; \omterm{def:in:how-pay-works:base}{salaire de base} médian
de 150\,000~\$~; le rapport ne dit pas combien ont déclaré un salaire. Carnegie Mellon MS in Computational Finance, promotion
2026 (diplômés en décembre 2025)~: 99 étudiants, 98 en recherche d'emploi, 98 avec une offre acceptée dans les trois mois, en comptant
les emplois de courte durée selon sa norme de déclaration~; \omterm{def:in:how-pay-works:base}{salaire de base} médian de 145\,000~\$ à partir de 89 étudiants ayant déclaré un salaire.
\end{dated}

\begin{method}[Ce qu'une médiane calculée sur une partie d'une promotion permet de dire]\label{met:in:education-pipelines:bounds}
Si $r$ diplômés sur $n$ déclarent un salaire, la médiane de l'ensemble des $n$ se situe entre les quantiles de la distribution déclarée
$\max\bigl(0,(n/2-(n-r))/r\bigr)$ et $\min\bigl(1,(n/2)/r\bigr)$, quel que soit le salaire des $n-r$ non-déclarants~: tous en dessous des
déclarants à une extrémité, tous au-dessus à l'autre.
\end{method}

Pour le second rapport, 89 sur 98 ont déclaré un salaire~: la médiane des 98 se situe entre les 44.9\up{e} et 55.1\up{e}~centiles du salaire
déclaré, un intervalle étroit. Pour le premier, le nombre de déclarants n'est pas indiqué, et aucune borne de ce type n'est possible. Aucun des deux
rapports ne dit ce que gagnaient les non-déclarants ni pourquoi ils n'ont pas déclaré, et un taux de 96~\% ou de 100~\% de «~placés~» dépend de ce
qui compte comme un emploi et de la date limite. Le coût et la durée des études doivent figurer dans la même comparaison.

\section{Doctorats et voie académique}

Les États-Unis ont décerné 1\,760 doctorats de recherche en physique, 1\,628 en informatique, 1\,165 en mathématiques et 454 en
statistique au cours de l'année se terminant en juin 2024 (la \cref{fig:in:education-pipelines:levels}), soit à peu près autant que les masters de
financial mathematics. Les doctorants entrent dans les rôles de recherche à un niveau supérieur à celui des nouveaux diplômés (chapitre~28), souvent
après un \omterm{def:in:career-paths-and-moves:entry}{stage} pendant la thèse. La voie académique fonctionne aussi dans l'autre sens~: des chercheurs quittent l'industrie pour des postes d'enseignant-chercheur,
et des enseignants-chercheurs font du conseil pour des entreprises ou les rejoignent~; aucun de ces flux n'est recensé par une source publique
que ce livre ait trouvée.

\begin{omfigure}
\begin{tikzpicture}
\begin{axis}[xbar,area legend,width=10.5cm,height=5.8cm,xmode=log,xmin=1,xmax=100000,log origin=infty,ymin=0.4,ymax=6.6,
  bar width=5pt,ytick=data,yticklabels={{physique},{informatique},{mathématiques},{statistique},{analyse financière},{mathématiques financières}},
  yticklabel style={font=\footnotesize},xticklabel style={font=\footnotesize},
  xlabel={\small diplômes décernés, année se terminant en juin 2024 (échelle logarithmique)},grid=major,grid style={gray!20},table/col sep=comma,
  legend style={legend cell align=left,font=\scriptsize,at={(0.5,-0.36)},anchor=north,legend columns=3,/tikz/every even column/.append style={column sep=6pt}}]
\addplot[fill=black!25,draw=none] table[x=bachelor,y=pos]{figdata/industry/29-education-pipelines/levels.csv};
\addplot[fill=omDef,draw=none] table[x=master,y=pos]{figdata/industry/29-education-pipelines/levels.csv};
\addplot[fill=omThm,draw=none] table[x=doctorate,y=pos]{figdata/industry/29-education-pipelines/levels.csv};
\legend{licence,masters,doctorat de recherche}
\end{axis}
\end{tikzpicture}

\omcaption{Diplômes américains dans six domaines par niveau, année se terminant en juin 2024 (IPEDS, premières filières). Données~:
\texttt{data/industry/ipeds\_completions.csv}, via \texttt{in\_edu.completions}.}\label{fig:in:education-pipelines:levels}
\end{omfigure}

\section{Concours et ce qu'ils signalent}

Les entreprises qui recrutent des diplômés en quantitatif lisent les résultats de concours comme une preuve de résolution de problèmes sous contrainte de temps. Le
plus connu est l'Olympiade internationale de mathématiques, réservée aux élèves «~âgés de moins de 20~ans et non inscrits dans un enseignement
postsecondaire~»~: des équipes nationales de six au plus, «~2~jours, 3~problèmes par jour, 4.5~heures chaque jour~», chaque problème
valant 7~points. Les concours universitaires de programmation et de mathématiques jouent un rôle analogue plus tard. Ce qu'un résultat signale est
étroit~: la rapidité et la profondeur sur des problèmes bien posés. La recherche demande des problèmes que personne n'a posés, des données qui mentent et de la patience
sur plusieurs mois~; le livre~18 montre comment les entretiens tentent de mesurer les deux.

\section{Apprendre en dehors d'un cursus}

Une grande partie de ce qu'utilise l'industrie s'apprend en dehors d'un diplôme~: la programmation en production, la statistique sur des données réelles, les
marchés eux-mêmes. Les voies sont ouvertes~: manuels et cours en ligne ouverts, bibliothèques open source (l'ingénierie publique du chapitre~20), concours
de données, et le projet en cours de cette série, dont le code accompagne chaque livre. Un candidat
sans diplôme spécialisé montre les mêmes preuves qu'un cursus~: un travail qui peut être lu et relancé.

\section{Tutoriel~: combien de quants le monde forme-t-il~?}\label{tut:in:education-pipelines:tut}

\begin{tutorial}
\textbf{Objectif.} Compter les diplômes américains dans les domaines de l'industrie, lire deux \omterm{def:in:education-pipelines:placement}{rapports d'insertion professionnelle} pour leur taux de réponse, et mettre
les diplômes spécialisés en regard des dépôts du chapitre~14.
\textbf{État final~:} \cref{fig:in:education-pipelines:masters,fig:in:education-pipelines:levels} et les bornes d'insertion.
\begin{tutsteps}
\item \textbf{Diplômes délivrés.} \texttt{firm.edupipe.read\_ipeds} parcourt le zip d'une année et totalise les premières filières par code
et par niveau (\cref{lst:in:education-pipelines:ipeds})~; \texttt{in\_ipeds\_derive.py} l'exécute sur quatre années.
\omcode{code/firm/edupipe/firm_edupipe.py}{23}{36}{Diplômes par code de programme et par niveau, à partir du fichier des diplômes d'une année.}\label{lst:in:education-pipelines:ipeds}
\item \textbf{Rapports d'insertion.} \texttt{PlacementReport} enregistre les dénominateurs~; \texttt{median\_quantile\_bounds} donne
la borne de la méthode (\cref{lst:in:education-pipelines:bounds}).
\omcode{code/firm/edupipe/firm_edupipe.py}{59}{66}{Les quantiles entre lesquels doit se situer la médiane d'une promotion.}\label{lst:in:education-pipelines:bounds}
\item \textbf{Le ratio.} Masters en deux codes financiers, rapportés aux dépôts de l'exercice 2025 dans les familles
de \omterm{def:in:quant-researcher:def}{chercheur quantitatif} et de développeur quant.
\end{tutsteps}
\end{tutorial}

Le résultat~: 5\,454 masters délivrés en 2024 contre 1\,466 dépôts pour l'exercice 2025, soit un ratio de 3.72. Les deux nombres
ne comptent pas la même chose (diplômes de toutes nationalités contre dépôts de visas d'un échantillon d'employeurs), et le ratio est une
échelle, non un taux d'insertion.

\textbf{Pour aller plus loin.} Ajouter les fichiers des autres années et les effectifs du Royaume-Uni lorsqu'ils sont sous licence ouverte~; ventiler les diplômes
par citoyenneté (l'enquête distingue les étudiants non résidents)~; recueillir davantage de rapports d'insertion avec leurs dénominateurs.

\section{Construction~: firm.edupipe}\label{bld:in:education-pipelines:edupipe}

\begin{build}
\bfield{Rôle} Compter les diplômes par domaine et par niveau à partir de l'enquête fédérale, et enregistrer les rapports d'insertion avec
leurs dénominateurs.
\bfield{Interface} \texttt{firm.edupipe}~: \texttt{LEVELS}~; \texttt{read\_ipeds(zip\_path, cips, levels)}~;
\texttt{PlacementReport(programme, cohort, students, seeking, accepted, reporting, median\_base, standard, note)} avec
\texttt{placement\_rate}, \texttt{reporting\_rate} et \texttt{median\_quantile\_bounds()}.
\bfield{Règles} Premières filières seulement~; le fichier est lu en flux~; un rapport qui n'indique pas son nombre de répondants ne donne aucune borne.
\bfield{Tests d'acceptation} \texttt{code/firm/edupipe/tests/}~: les deuxièmes filières et les autres codes sont exclus, et les niveaux à zéros de tête
sont lus~; 89 rapports sur 98 bornent la médiane entre le 44.9\up{e}~centile et le 55.1\up{e}~centile publiés.
\bfield{Pour aller plus loin} Les comptages HESA~; la part des non-résidents parmi les diplômés~; un analyseur des rapports d'insertion.
\end{build}

\begin{omsources}
\item National Center for Education Statistics, IPEDS completions et Classification of Instructional Programs 2020.
\item UC Berkeley MFE et Carnegie Mellon MSCF employment reports.
\item International Mathematical Olympiad, general information.
\item Tableaux du chapitre~14 issus des fichiers LCA du Department of Labor.
\end{omsources}

\section{Exercices}

\begin{exercise}[$\star$]\label{exo:in:education-pipelines:1}
Déterminer de quel facteur les masters en mathématiques financières ont augmenté entre 2010 et 2024.
\end{exercise}

\begin{exercise}[$\star$]\label{exo:in:education-pipelines:2}
Déterminer quelle part de la promotion de Berkeley à la recherche d'un emploi a accepté une offre.
\end{exercise}

\begin{exercise}[$\star$]\label{exo:in:education-pipelines:3}
Déterminer le nombre maximal de points qu'un candidat à l'OIM peut obtenir.
\end{exercise}

\begin{exercise}[$\star\star$]\label{exo:in:education-pipelines:4}
Si 60 des 98 rapports ont publié les salaires dans le second rapport, déterminer entre quels centiles publiés la médiane de la promotion doit se situer.
\end{exercise}

\begin{exercise}[$\star\star$]\label{exo:in:education-pipelines:5}
Expliquer pourquoi le nombre de diplômes en mathématiques financières constitue un plancher pour les diplômes financiers spécialisés, et en quoi il pourrait
surestimer la tendance.
\end{exercise}

\begin{exercise}[$\star\star$]\label{exo:in:education-pipelines:6}
Comparer les doctorats de recherche dans les quatre grands domaines généraux avec les masters financiers de 2024.
\end{exercise}

\begin{exercise}[$\star\star\star$]\label{exo:in:education-pipelines:7}
\emph{Programmation.} À l'aide de \texttt{read\_ipeds}, calculer la croissance des masters en informatique entre 2019 et 2024,
et la comparer à celle des mathématiques financières.
\end{exercise}

\begin{exercise}[$\star\star\star$]\label{exo:in:education-pipelines:8}
\emph{Trouver la faille.} «~Le programme A affiche une médiane de 150\,000~\$ et le programme B de 145\,000~\$, donc A place
mieux ses diplômés~».
\end{exercise}

\section{Problème~: combien de quants le monde forme-t-il~?}

\begin{problem}[{Problème du week-end --- combien de quants le monde forme-t-il~?}]\label{pb:in:education-pipelines:1}
Une université envisage de lancer un master en ingénierie financière et demande quelle est la taille du marché de ses diplômés.

\medskip\noindent\textbf{Partie I --- Les codes.}
\begin{enumerate}
\item Définir la \omterm{def:in:education-pipelines:cip}{Classification of Instructional Programs} (CIP).
\item Quels codes couvrent les mathématiques financières et l'analyse financière, et comment sont-ils définis~?
\item Que compte l'enquête sur les diplômes, et sur quelle période~?
\item Pourquoi les premières filières comptent-elles~?
\item Pourquoi le comptage constitue-t-il un plancher~?
\end{enumerate}

\medskip\noindent\textbf{Partie II --- Les chiffres.}
\begin{enumerate}[resume]
\item Donner le nombre de masters en mathématiques financières pour les quatre années.
\item Donner les diplômes délivrés en 2024 par niveau dans les six domaines.
\item Quelle est l'importance des diplômes spécialisés à côté de l'informatique~?
\item Donner les doctorats de recherche dans les quatre grands domaines.
\item Quels métiers mobilisent quels diplômes~?
\end{enumerate}

\medskip\noindent\textbf{Partie III --- Insertion.}
\begin{enumerate}[resume]
\item Définir un \omterm{def:in:education-pipelines:placement}{rapport d'insertion professionnelle}.
\item Donner les dénominateurs et les médianes des deux rapports.
\item Énoncer la borne sur la médiane d'une promotion et l'appliquer.
\item Que ne dit ni l'un ni l'autre rapport~?
\item Que signale le résultat d'un concours, et que ne signale-t-il pas~?
\end{enumerate}

\medskip\noindent\textbf{Partie IV --- Le verdict.}
\begin{enumerate}[resume]
\item Énoncer le \emph{résultat nommé}~: les diplômes de master annuels délivrés aux États-Unis dans les codes de filières financières, et leur rapport
au nombre annuel de déclarations LCA portant sur des intitulés quantitatifs chez des employeurs de la finance.
\item Pourquoi ce rapport n'est-il pas un taux d'insertion~?
\item À quoi les diplômés de l'université seraient-ils en concurrence~?
\item Que devrait publier son \omterm{def:in:education-pipelines:placement}{rapport d'insertion professionnelle}~?
\item En deux phrases, répondre à l'université.
\end{enumerate}
\end{problem}

\section{Questions d'entretien}

\begin{interviewq}[$\star$ \iqroles{chercheur}]\label{iq:in:education-pipelines:1}
Qu'est-ce que votre thèse vous a appris qu'un cours ne pouvait pas vous enseigner~?
\end{interviewq}

\begin{interviewq}[$\star$ \iqroles{chercheur, développeur}]\label{iq:in:education-pipelines:2}
Montrez-nous un travail réalisé en dehors de vos cours et expliquez-nous une décision qu'il contient.
\end{interviewq}

\begin{interviewq}[$\star\star$ \iqroles{chercheur}]\label{iq:in:education-pipelines:3}
Un rapport affirme que 100~\% d'une promotion a été insérée. Que demanderiez-vous avant d'y croire~?
\end{interviewq}

\begin{interviewq}[$\star\star$ \iqroles{chercheur}]\label{iq:in:education-pipelines:4}
Comment estimeriez-vous le nombre de personnes qui deviennent \omterm{def:in:quant-researcher:def}{chercheurs quantitatifs} chaque année~?
\end{interviewq}

\begin{interviewq}[$\star\star$ \iqroles{développeur}]\label{iq:in:education-pipelines:5}
Vous vous êtes formé seul au C++. Comment nous montreriez-vous que vous savez écrire du code de production~?
\end{interviewq}

\begin{interviewq}[$\star\star\star$ \iqroles{chercheur}]\label{iq:in:education-pipelines:6}
Une série de comptages fait un saut lorsqu'un nouveau code de classification est introduit. Comment distinguez-vous un vrai changement d'un recodage~?
\end{interviewq}
